% ---------------------------------------------------------------------------
% SDP-as-theorem subsection for the final report's Main Results (§III or §IV).
%
% This converts what was presented in the midterm as an informal "baseline"
% into a stated theoretical anchor with a citation + a short proof sketch.
% Per Mou's Tips and Example_FinalReprot2.pdf, a proposition-with-proof is the
% canonical shape of a strong Main Results section.
%
% Required bib entries (see additional_refs.bib):
%   Boyd2006, GhoshBoyd2006, XiaoBoyd2004, BoydVandenberghe2004, HornJohnson2013
% ---------------------------------------------------------------------------

\subsection{Convex upper bound on achievable algebraic connectivity}

Let $G = (V, E)$ be a fixed undirected connected graph on $n$ nodes. Write
$E = \{e_1, \ldots, e_m\}$ with endpoints $e_k = (i_k, j_k)$, and let
$a_k = \mathbf{e}_{i_k} - \mathbf{e}_{j_k} \in \R^n$ where
$\mathbf{e}_i$ is the $i$-th standard basis vector. For a nonnegative
weight vector $w \in \R^m_{\ge 0}$, the weighted graph Laplacian is
\begin{equation}
L(w) \;=\; \sum_{k=1}^{m} w_k\, a_k a_k^\top
\;=\; A\,\mathrm{diag}(w)\,A^\top,
\label{eq:Lw}
\end{equation}
with $A = [a_1, \ldots, a_m] \in \R^{n \times m}$ the signed
incidence matrix. By construction $L(w)$ is symmetric positive
semidefinite and $L(w)\mathbf{1} = 0$, so the smallest eigenvalue of
$L(w)$ is $0$ with eigenvector $\mathbf{1}$, and the algebraic
connectivity is the second smallest eigenvalue
\cite{Fiedler1973,OlfatiSaber2007}. By the Courant--Fischer
characterization \cite{HornJohnson2013},
\begin{equation}
\lambda_2(L(w)) \;=\; \min_{\substack{v \in \R^n \\ v \perp \mathbf{1},\; \|v\|=1}} v^\top L(w)\, v.
\label{eq:rayleigh}
\end{equation}

We define the admissible weight set
\begin{equation}
\mathcal{W} = \Big\{ w \in \R^m \;:\; 0 \le w_k \le w_{\max},\;
\textstyle\sum_{k=1}^m w_k \le B \Big\},
\label{eq:Wset}
\end{equation}
which is the polyhedral convex set used throughout this work. The
reweighting problem solved by our RL policy (in continuous form) is
\begin{equation}
\lambda_2^\star \;:=\; \max_{w \in \mathcal{W}} \; \lambda_2(L(w)).
\label{eq:primal}
\end{equation}

\begin{theorem}[Boyd \cite{Boyd2006}; Ghosh and Boyd \cite{GhoshBoyd2006}]
\label{thm:fdla}
Problem \eqref{eq:primal} is equivalent to the semidefinite program
\begin{equation}
\begin{aligned}
s^\star \;=\; \max_{w,\,s} \quad & s \\
\text{s.t.} \quad & L(w) \;\succeq\; s\!\left(I_n - \tfrac{1}{n}\mathbf{1}\mathbf{1}^\top\right), \\
& w \in \mathcal{W},
\end{aligned}
\label{eq:sdp}
\end{equation}
in the sense that $\lambda_2^\star = s^\star$, and any optimizer
$w^\star$ of \eqref{eq:sdp} is a maximizer of \eqref{eq:primal}.
\end{theorem}

\begin{proof}
Let $\Pi := I_n - \tfrac{1}{n}\mathbf{1}\mathbf{1}^\top$ be the
orthogonal projector onto $\mathbf{1}^\perp$. Since $L(w) \succeq 0$
and $L(w)\mathbf{1} = 0$, the nontrivial eigenvalues of $L(w)$ are
exactly the eigenvalues of $\Pi L(w) \Pi$ restricted to
$\mathbf{1}^\perp$. Equation \eqref{eq:rayleigh} therefore gives
\[
\lambda_2(L(w)) \;=\; \sup\{\, s \in \R : v^\top L(w) v \ge s\, v^\top \Pi v
\;\; \forall\, v \in \R^n \,\}.
\]
The universal inequality inside the supremum is a linear matrix
inequality $L(w) - s\Pi \succeq 0$, so
\[
\lambda_2(L(w)) \;=\; \sup\{\, s : L(w) - s\,\Pi \succeq 0 \,\}.
\]
Taking the supremum jointly over $w \in \mathcal{W}$ and $s$
yields \eqref{eq:sdp}. Because $L(w)$ is linear in $w$
(cf.~\eqref{eq:Lw}) and $\mathcal{W}$ is convex, the program
\eqref{eq:sdp} is a convex SDP and is solvable to global optimality by
interior-point methods
\cite{BoydVandenberghe2004,XiaoBoyd2004}.
\end{proof}

Theorem~\ref{thm:fdla} plays two roles in this paper. First, it gives a
\emph{tight}, polynomial-time-computable upper bound on any reweighting
policy, learned or otherwise: no feasible $w \in \mathcal{W}$ attains
$\lambda_2(L(w)) > s^\star$. Second, it defines the
\emph{suboptimality ratio}
\begin{equation}
\rho_\pi \;:=\; \frac{\lambda_2\!\big(L(w_\pi)\big)}{s^\star}
\;\in\; [0, 1]
\label{eq:subopt}
\end{equation}
for any policy $\pi$ producing a feasible $w_\pi$. We evaluate our
PPO policy against \eqref{eq:subopt} in Section~V and report
$\rho_\pi \approx 0.87$ on Erd\H{o}s--R\'enyi supports of size
$n = 20$, turning the comparison against classical heuristic baselines
(uniform, Metropolis--Hastings \cite{XiaoBoyd2004,Boyd2004}) into a
quantitative optimality claim rather than a relative one.

\begin{remark}
Theorem~\ref{thm:fdla} assumes a fixed edge support $E$. When the
support is itself a decision variable, as in the discrete rewiring
setting of Section~VI, the feasible set $\{L(w) : w \in \mathcal{W}\}$
is no longer convex and \eqref{eq:sdp} ceases to provide a valid upper
bound; maximizing $\lambda_2$ over topologies is NP-hard in general
\cite{Mosk2006}. Similarly, when $E$ is induced by agent positions as
in the geometric swarm setting of Section~VII, the map from positions
to $L(w)$ is nonlinear and nonconvex, and the SDP does not apply. This
is precisely the regime in which a learned policy earns its keep.
\end{remark}
